%% kapa: exam=1 round=28 year=2017 period=2 form=A subject=appraisal_law subject_ko=감정평가관계법규 seq=5 number=5
%% source: Q-Net 감정평가사 기출문제 (cst00309) artlSeq=5199603 file=extracted/1008/5199603/2017년 제28회 감정평가사 1차시험 문제지/2교시 A형.pdf
%% conversion: kordoc 4.24.51 -> tools/gen1.py (auto); review=proofread-sonnet; flags=none
\begin{question}{5}
국토의 계획 및 이용에 관한 법령상 용도구역에 관한 설명이다. 다음 ( )안에 알맞은 것은?
\begin{framedtext}
시ㆍ도지사 또는 대도시 시장은 도시의 자연환경 및 경관을 보호하고 도시민에게 건전한 여가ㆍ휴식공간을 제공하기 위하여 도시지역 안에서 식생이 양호한 산지(山地)의 개발을 제한할 필요가 있다고 인정하면 ( )의 지정 또는 변경을 도시ㆍ군관리계획으로 결정할 수 있다.
\end{framedtext}
\begin{choices}
  \item 시가화조정구역
  \item 개발제한구역
  \item 수산자원보호구역
  \item 입지규제최소구역
  \item 도시자연공원구역
\end{choices}
\end{question}
\answer{5}
